% ============================================================================
%  PRECÁLCULO — Sesión 9
%  Semana 9: 5 al 10 de octubre de 2026
%  2.3 Racionalización y sus procesos
%  Duración prevista: 3 horas
% ============================================================================
\documentclass[aspectratio=169,10pt]{beamer}

\usepackage[utf8]{inputenc}
\IfFileExists{spanish.ldf}{\usepackage[spanish,es-noshorthands]{babel}}{}
\usepackage{amsmath,amssymb}
\usepackage{booktabs}
\usepackage{tikz}

\usetheme{Uniminuto}

\uniasignatura{Precálculo}
\uniprofesor{Andrés Felipe Puerta Velez}
\unicorreo{andres.puerta-v@uniminuto.edu.co}
\uninrc{MATE1161}
\unisemana{Sesión 9 — Semana 9 · 5 al 10 de octubre de 2026}

\newcommand{\listatiny}{\setbeamerfont{itemize/enumerate body}{size=\tiny}\tiny}

\begin{document}

\portada

\begin{frame}{Orden del día}
  \begin{enumerate}
    \item Llamado a lista y revisión de los ejercicios propuestos
    \item Repaso: dónde quedamos
    \item \azul{Tema 2.3: por qué y cómo se racionaliza}
    \item Caso monomio: raíz cuadrada y raíz de índice $n$
    \item Caso binomio: el conjugado
    \item Racionalizar el numerador; taller y cierre
  \end{enumerate}

  \vspace{0.25cm}
  {\footnotesize\color{uniGris} Sesión de 3 horas. Todo lo de hoy se apoya en
  una sola idea de la sesión 8: \destacar{multiplicar por 1 escrito de forma
  conveniente}.}
\end{frame}

\begin{frame}{Repaso: dónde quedamos}
  \footnotesize
  En la sesión 8 simplificamos radicales, sumamos semejantes y operamos
  productos y cocientes de raíces.

  \vspace{0.15cm}
  \begin{alertblock}{Las dos herramientas que hoy se usan todo el tiempo}
    \scriptsize
    \[ \left(\sqrt{a}\right)^{2} = a \qquad\qquad
       (a+b)(a-b) = a^{2}-b^{2} \]
    La primera es la definición de raíz. La segunda es la
    \azul{diferencia de cuadrados} que apareció en el ejercicio 3 de la sesión
    pasada: $(\sqrt5+\sqrt3)(\sqrt5-\sqrt3) = 5-3 = 2$.
  \end{alertblock}

  \vspace{0.15cm}
  \begin{block}{La pregunta de hoy}
    \footnotesize
    \destacar{¿Cómo se quita una raíz del denominador sin cambiar el valor de
    la fracción?}
  \end{block}
\end{frame}

\seccion{2.3 Racionalización}

\begin{frame}{Por qué racionalizar}
  \scriptsize
  \begin{columns}[T,onlytextwidth]
    \begin{column}{0.48\textwidth}
      \begin{block}{Las tres razones}
        \tiny
        \textbf{1. Forma canónica.} Un resultado se entrega sin raíces en el
        denominador; es la convención del curso y de los textos.

        \vspace{0.06cm}
        \textbf{2. Comparar y sumar.} Con denominador entero es inmediato
        buscar denominador común.

        \vspace{0.06cm}
        \textbf{3. Cálculo numérico.} Dividir entre un irracional propaga más
        error que multiplicar por él.
      \end{block}
    \end{column}
    \begin{column}{0.48\textwidth}
      \begin{exampleblock}{Un ejemplo del punto 3}
        \tiny
        \[ \frac{1}{\sqrt{2}} \quad\text{frente a}\quad \frac{\sqrt{2}}{2} \]
        Con $\sqrt2 \approx 1{,}414$:
        \[ 1 \div 1{,}414 = 0{,}7072\dots \]
        \[ 1{,}414 \div 2 = 0{,}707 \]
        La segunda es \azul{una división entre 2}: exacta y sin arrastre de
        error.
      \end{exampleblock}
    \end{column}
  \end{columns}

  \vspace{0.1cm}
  \begin{alertblock}{La única regla del tema}
    \scriptsize
    \destacar{Se multiplica arriba y abajo por lo mismo.} Eso es multiplicar
    por $1$, así que \azul{el valor de la fracción no cambia}: solo cambia su
    aspecto. Todo el trabajo consiste en elegir bien por qué multiplicar.
  \end{alertblock}
\end{frame}

\begin{frame}{Caso 1 — un solo radical cuadrado en el denominador}
  \scriptsize
  \begin{block}{Se multiplica por esa misma raíz}
    \tiny
    \[ \frac{a}{\sqrt{b}}\cdot\frac{\sqrt{b}}{\sqrt{b}}
       = \frac{a\sqrt{b}}{b} \]
  \end{block}

  \vspace{0.08cm}
  \begin{columns}[T,onlytextwidth]
    \begin{column}{0.48\textwidth}
      \begin{exampleblock}{Ejemplo 1}
        \tiny
        \[ \frac{3}{\sqrt{5}} = \frac{3}{\sqrt5}\cdot\frac{\sqrt5}{\sqrt5}
           = \destacar{\frac{3\sqrt{5}}{5}} \]
      \end{exampleblock}
      \begin{exampleblock}{Ejemplo 2}
        \tiny
        \[ \frac{7}{2\sqrt{3}} = \frac{7\sqrt3}{2\cdot 3}
           = \destacar{\frac{7\sqrt{3}}{6}} \]
        {\color{uniGris} El 2 no se toca: solo estorba la raíz.}
      \end{exampleblock}
    \end{column}
    \begin{column}{0.48\textwidth}
      \begin{exampleblock}{Ejemplo 3 — simplificar primero}
        \tiny
        \[ \frac{8}{\sqrt{2}} = \frac{8\sqrt2}{2} = \destacar{4\sqrt{2}} \]
      \end{exampleblock}
      \begin{alertblock}{Cuidado}
        \tiny
        Después de racionalizar, \destacar{casi siempre se puede simplificar}.
        Dejar $\tfrac{8\sqrt2}{2}$ como respuesta final es un ejercicio a
        medias.
      \end{alertblock}
    \end{column}
  \end{columns}
\end{frame}

\begin{frame}{Caso 2 — raíz de índice $n$}
  \scriptsize
  \begin{block}{Se completa el exponente hasta el índice}
    \tiny
    $\dfrac{a}{\sqrt[n]{b^{m}}}\cdot\dfrac{\sqrt[n]{b^{\,n-m}}}
    {\sqrt[n]{b^{\,n-m}}} = \dfrac{a\sqrt[n]{b^{\,n-m}}}{b}$ \quad
    La idea: dentro de la raíz hay $b^{m}$ y hace falta llegar a $b^{n}$ para
    que la raíz \azul{desaparezca}.
  \end{block}

  \vspace{0.05cm}
  \begin{columns}[T,onlytextwidth]
    \begin{column}{0.48\textwidth}
      \begin{exampleblock}{Ejemplo 1}
        \tiny
        Falta $2^{2}$ para completar $2^{3}$:
        \[ \frac{5}{\sqrt[3]{2}}\cdot\frac{\sqrt[3]{2^{2}}}{\sqrt[3]{2^{2}}}
           = \frac{5\sqrt[3]{4}}{\sqrt[3]{8}}
           = \destacar{\frac{5\sqrt[3]{4}}{2}} \]
      \end{exampleblock}
    \end{column}
    \begin{column}{0.48\textwidth}
      \begin{exampleblock}{Ejemplo 2}
        \tiny
        Como $\sqrt[3]{9}=\sqrt[3]{3^{2}}$, falta un $3$:
        \[ \frac{2}{\sqrt[3]{9}}\cdot\frac{\sqrt[3]{3}}{\sqrt[3]{3}}
           = \frac{2\sqrt[3]{3}}{\sqrt[3]{27}}
           = \destacar{\frac{2\sqrt[3]{3}}{3}} \]
      \end{exampleblock}
    \end{column}
  \end{columns}

  \vspace{0.05cm}
  \begin{alertblock}{El error que se comete siempre}
    \tiny
    Multiplicar por $\sqrt[3]{2}$ en vez de $\sqrt[3]{2^{2}}$: eso deja
    $\sqrt[3]{4}$ abajo y \destacar{la raíz no desapareció}.
  \end{alertblock}
\end{frame}

\begin{frame}{Caso 3 — binomio con raíces: el conjugado}
  \scriptsize
  {\tiny El conjugado de $a+b$ es $a-b$, y su producto es $a^{2}-b^{2}$:
  \destacar{una diferencia de cuadrados que elimina las raíces}.}

  \vspace{0.08cm}
  \begin{columns}[T,onlytextwidth]
    \begin{column}{0.50\textwidth}
      \begin{exampleblock}{Ejemplo guiado}
        \tiny
        \[ \frac{2}{\sqrt{7}-\sqrt{3}}
           \cdot\frac{\sqrt{7}+\sqrt{3}}{\sqrt{7}+\sqrt{3}} \]
        \[ = \frac{2(\sqrt7+\sqrt3)}{(\sqrt7)^{2}-(\sqrt3)^{2}}
           = \frac{2(\sqrt7+\sqrt3)}{7-3} \]
        \[ = \frac{2(\sqrt7+\sqrt3)}{4}
           = \destacar{\frac{\sqrt{7}+\sqrt{3}}{2}} \]
      \end{exampleblock}
    \end{column}
    \begin{column}{0.46\textwidth}
      \begin{exampleblock}{Otro caso}
        \tiny
        \[ \frac{6}{\sqrt{5}+2}\cdot\frac{\sqrt5-2}{\sqrt5-2}
           = \frac{6(\sqrt5-2)}{5-4} \]
        \[ = \destacar{6\sqrt{5}-12} \]
        {\color{uniGris} Aquí el denominador quedó en $1$ y desapareció.}
      \end{exampleblock}
    \end{column}
  \end{columns}

  \vspace{0.05cm}
  \begin{alertblock}{Lo único que cambia es el signo del medio}
    \tiny
    Conjugado de $\sqrt5+2$ es $\sqrt5-2$. \quad De $3-\sqrt2$ es $3+\sqrt2$.
    \destacar{No se cambia el orden de los términos ni se tocan las raíces:}
    solo el signo que los separa.
  \end{alertblock}
\end{frame}

\begin{frame}{Racionalizar el numerador}
  \scriptsize
  \begin{columns}[T,onlytextwidth]
    \begin{column}{0.50\textwidth}
      \begin{block}{El mismo procedimiento, al revés}
        \tiny
        A veces conviene dejar la raíz abajo y quitarla de arriba:
        \[ \frac{\sqrt{x+h}-\sqrt{x}}{h}
           \cdot\frac{\sqrt{x+h}+\sqrt{x}}{\sqrt{x+h}+\sqrt{x}} \]
        \[ = \frac{(x+h)-x}{h\left(\sqrt{x+h}+\sqrt{x}\right)}
           = \destacar{\frac{1}{\sqrt{x+h}+\sqrt{x}}} \]
      \end{block}
    \end{column}
    \begin{column}{0.46\textwidth}
      \begin{alertblock}{Por qué se lo muestro hoy}
        \tiny
        Esa expresión es el \destacar{cociente incremental} con el que se
        define la derivada en Cálculo Diferencial.

        \vspace{0.08cm}
        Tal como está, al hacer $h \to 0$ queda $\tfrac{0}{0}$. Racionalizada,
        \azul{el problema desaparece}.

        \vspace{0.08cm}
        No hay que resolverlo hoy: solo ver que \destacar{esta técnica no se
        queda en Precálculo}.
      \end{alertblock}
    \end{column}
  \end{columns}

  \vspace{0.08cm}
  {\tiny\color{uniGris} Regla práctica: \destacar{se racionaliza donde estorbe}.
  En una fracción normal, el denominador; en un límite, casi siempre el
  numerador.}
\end{frame}

\begin{frame}{Ejercicio 1 — en clase}
  \begin{alertblock}{Individual (12 minutos)}
    \scriptsize
    Racionalice y simplifique:

    \vspace{0.3cm}
    \centering
    \begin{tabular}{@{}llll@{}}
      1. \ $\dfrac{8}{\sqrt{2}}$ &
      2. \ $\dfrac{\sqrt{3}}{\sqrt{7}}$ &
      3. \ $\dfrac{5}{2\sqrt{5}}$ &
      4. \ $\dfrac{9}{\sqrt{3}}$ \\[1.3em]
      5. \ $\dfrac{1}{\sqrt[3]{3}}$ &
      6. \ $\dfrac{2}{\sqrt[3]{9}}$ &
      7. \ $\dfrac{4}{\sqrt[4]{8}}$ &
      8. \ $\dfrac{\sqrt{6}}{\sqrt{10}}$ \\
    \end{tabular}
  \end{alertblock}
\end{frame}

\begin{frame}{Ejercicio 1 — solución}
  \scriptsize
  \begin{center}
  \begin{tabular}{@{}clc@{\hspace{0.9cm}}clc@{}}
    \toprule
    \textbf{\#} & \textbf{Se multiplica por} & \textbf{Resultado} &
    \textbf{\#} & \textbf{Se multiplica por} & \textbf{Resultado} \\
    \midrule
    1 & $\sqrt2$ & $4\sqrt{2}$ &
    5 & $\sqrt[3]{3^{2}}$ & $\dfrac{\sqrt[3]{9}}{3}$ \\[0.4em]
    2 & $\sqrt7$ & $\dfrac{\sqrt{21}}{7}$ &
    6 & $\sqrt[3]{3}$ & $\dfrac{2\sqrt[3]{3}}{3}$ \\[0.4em]
    3 & $\sqrt5$ & $\dfrac{\sqrt{5}}{2}$ &
    7 & $\sqrt[4]{2}$ & $2\sqrt[4]{2}$ \\[0.4em]
    4 & $\sqrt3$ & $3\sqrt{3}$ &
    8 & $\sqrt{10}$ & $\dfrac{\sqrt{15}}{5}$ \\
    \bottomrule
  \end{tabular}
  \end{center}

  \vspace{0.1cm}
  \begin{alertblock}{Los dos que exigen atención}
    \tiny
    \textbf{3.} $\tfrac{5\sqrt5}{2\cdot 5} = \tfrac{\sqrt5}{2}$: hay que
    \destacar{simplificar el 5}. \quad
    \textbf{7.} $\sqrt[4]{8}=\sqrt[4]{2^3}$, falta un $2$:
    $\tfrac{4\sqrt[4]{2}}{2} = 2\sqrt[4]{2}$.
  \end{alertblock}
\end{frame}

\begin{frame}{Ejercicio 2 — en clase}
  \begin{alertblock}{En parejas (18 minutos)}
    \scriptsize
    Racionalice con el conjugado y simplifique:

    \vspace{0.3cm}
    \begin{columns}[T,onlytextwidth]
      \begin{column}{0.48\textwidth}
        \begin{enumerate}\setlength{\itemsep}{0.24cm}
          \item $\dfrac{12}{\sqrt{5}-1}$
          \item $\dfrac{1}{2+\sqrt{3}}$
          \item $\dfrac{14}{3+\sqrt{2}}$
        \end{enumerate}
      \end{column}
      \begin{column}{0.48\textwidth}
        \begin{enumerate}\setlength{\itemsep}{0.24cm}
          \setcounter{enumi}{3}
          \item $\dfrac{\sqrt{7}-\sqrt{5}}{\sqrt{7}+\sqrt{5}}$
          \item $\dfrac{1}{\sqrt{5}-\sqrt{2}}$
          \item $\dfrac{1}{\sqrt{x}+2}$ \ (con $x \ge 0$, $x\neq 4$)
        \end{enumerate}
      \end{column}
    \end{columns}
  \end{alertblock}
\end{frame}

\begin{frame}{Ejercicio 2 — solución}
  \tiny
  \begin{columns}[T,onlytextwidth]
    \begin{column}{0.49\textwidth}
      \begin{block}{Puntos 1 a 3}
        \tiny
        \begin{align*}
          1)\;& \frac{12(\sqrt5+1)}{5-1} = \frac{12(\sqrt5+1)}{4}
                = \destacar{3\sqrt{5}+3} \\[0.3em]
          2)\;& \frac{2-\sqrt3}{4-3} = \destacar{2-\sqrt{3}} \\[0.3em]
          3)\;& \frac{14(3-\sqrt2)}{9-2} = \frac{14(3-\sqrt2)}{7}
                = \destacar{6-2\sqrt{2}}
        \end{align*}
      \end{block}
    \end{column}
    \begin{column}{0.49\textwidth}
      \begin{block}{Puntos 4 a 6}
        \tiny
        \begin{align*}
          4)\;& \frac{(\sqrt7-\sqrt5)^{2}}{7-5}
                = \frac{12-2\sqrt{35}}{2} = \destacar{6-\sqrt{35}} \\[0.3em]
          5)\;& \frac{\sqrt5+\sqrt2}{5-2}
                = \destacar{\frac{\sqrt{5}+\sqrt{2}}{3}} \\[0.3em]
          6)\;& \destacar{\frac{\sqrt{x}-2}{x-4}}
        \end{align*}
      \end{block}
    \end{column}
  \end{columns}

  \vspace{0.08cm}
  \begin{alertblock}{El punto 4, con detalle}
    \tiny
    Al multiplicar arriba y abajo por $\sqrt7-\sqrt5$, el numerador queda
    $(\sqrt7-\sqrt5)^2 = 7 - 2\sqrt{35} + 5 = 12-2\sqrt{35}$. \destacar{Ahí
    sí hay que desarrollar el cuadrado del binomio}, no es una diferencia de
    cuadrados. \azul{El conjugado va abajo; arriba se multiplica normal.}
  \end{alertblock}
\end{frame}

\begin{frame}{Ejercicio 3 — en clase (taller principal)}
  \begin{alertblock}{En grupos de 3 (20 minutos)}
    \tiny
    \setbeamerfont{itemize/enumerate body}{size=\tiny}
    \begin{enumerate}\setlength{\itemsep}{0.05cm}
      \item Racionalice: \ $\dfrac{3}{\sqrt{11}+3}$
      \item Racionalice: \ $\dfrac{\sqrt{2}}{\sqrt{6}-2}$
      \item Racionalice: \ $\dfrac{10}{\sqrt{6}+1}$
      \item Sume, racionalizando primero:
            \ $\dfrac{1}{\sqrt{2}} + \dfrac{1}{\sqrt{8}}$
      \item Verifique numéricamente que $\dfrac{1}{2+\sqrt3}$ y $2-\sqrt3$ son
            el mismo número (use $\sqrt3\approx 1{,}732$).
      \item Encuentre el error: \\
            \hspace*{0.3cm}
            $\dfrac{4}{\sqrt{3}} = \dfrac{4\sqrt3}{\sqrt3\cdot\sqrt3}
            = \dfrac{4\sqrt3}{9}$
    \end{enumerate}
  \end{alertblock}
\end{frame}

\begin{frame}{Ejercicio 3 — solución}
  \scriptsize
  \begin{columns}[T,onlytextwidth]
    \begin{column}{0.48\textwidth}
      \begin{block}{Puntos 1 a 4}
        \tiny
        \textbf{1.} $\dfrac{3(\sqrt{11}-3)}{11-9}
        = \destacar{\dfrac{3\sqrt{11}-9}{2}}$

        \vspace{0.06cm}
        \textbf{2.} $\dfrac{\sqrt2(\sqrt6+2)}{6-4}
        = \dfrac{\sqrt{12}+2\sqrt2}{2} = \destacar{\sqrt{3}+\sqrt{2}}$

        \vspace{0.06cm}
        \textbf{3.} $\dfrac{10(\sqrt6-1)}{6-1} = \destacar{2\sqrt{6}-2}$

        \vspace{0.06cm}
        \textbf{4.} $\tfrac{\sqrt2}{2} + \tfrac{\sqrt2}{4}
        = \destacar{\tfrac{3\sqrt{2}}{4}}$
      \end{block}
    \end{column}
    \begin{column}{0.50\textwidth}
      \begin{alertblock}{Puntos 5 y 6}
        \tiny
        \textbf{5.} $\tfrac{1}{2+1{,}732} = \tfrac{1}{3{,}732}
        = 0{,}268$ \ y \ $2-1{,}732 = 0{,}268$ \checkmark
        \azul{Racionalizar no cambia el valor, solo la forma.}

        \vspace{0.08cm}
        \textbf{6.} El error está en el denominador:
        $\sqrt3\cdot\sqrt3 = \destacar{3}$, no $9$. Lo correcto es
        $\tfrac{4\sqrt3}{3}$.
      \end{alertblock}
    \end{column}
  \end{columns}

  \vspace{0.05cm}
  {\tiny\color{uniGris} En el \destacar{2}, $\sqrt2\cdot\sqrt6 = \sqrt{12}
  = 2\sqrt3$, y al dividir entre 2 queda $\sqrt3$. \azul{Simplificar al final
  es parte del ejercicio}, no un extra.}
\end{frame}

\begin{frame}{Errores frecuentes en esta sesión}
  \scriptsize
  \begin{center}
  \begin{tabular}{@{}p{5.2cm}p{6.6cm}@{}}
    \toprule
    \textbf{Error} & \textbf{Corrección} \\
    \midrule
    $\sqrt{a}\cdot\sqrt{a} = a^{2}$ &
      Es $a$: la raíz y el cuadrado se cancelan \\
    \addlinespace
    Multiplicar solo el denominador &
      Se multiplica \azul{arriba y abajo}, o la fracción cambia \\
    \addlinespace
    Con $\sqrt[3]{2}$, multiplicar por $\sqrt[3]{2}$ &
      Hay que multiplicar por $\sqrt[3]{2^{2}}$ para completar el índice \\
    \addlinespace
    Cambiar el signo de los dos términos del conjugado &
      Solo cambia el signo \azul{del medio} \\
    \addlinespace
    Aplicar diferencia de cuadrados en el numerador &
      Arriba se multiplica normal; si hay binomio, se desarrolla \\
    \addlinespace
    Dejar el resultado sin simplificar &
      Casi siempre queda un factor común por cancelar \\
    \bottomrule
  \end{tabular}
  \end{center}
\end{frame}

\begin{frame}{Cierre de la sesión}
  \begin{block}{Lo que debe quedar claro hoy}
    \scriptsize
    \begin{itemize}\setlength{\itemsep}{0.02cm}
      \item Racionalizar es \azul{multiplicar por 1} escrito de forma
            conveniente: el valor no cambia.
      \item Monomio con $\sqrt{\ }$: se multiplica por esa misma raíz.
      \item Raíz de índice $n$: se \azul{completa el exponente} hasta el
            índice.
      \item Binomio: se usa el \azul{conjugado}, que produce una diferencia de
            cuadrados.
      \item A veces se racionaliza el \destacar{numerador}: es lo que se hará
            en Cálculo.
      \item Después de racionalizar, \destacar{siempre} se revisa si se puede
            simplificar.
    \end{itemize}
  \end{block}

  \vspace{0.08cm}
  \begin{alertblock}{Trabajo independiente}
    \tiny
    Resolver los 20 ejercicios de la lámina siguiente. En la sesión 10 empiezan
    los \destacar{logaritmos}: la operación que despeja el exponente y que
    explica por qué la búsqueda binaria es tan rápida.
  \end{alertblock}
\end{frame}

\begin{frame}{Ejercicios propuestos}
  \vspace{-0.1cm}
  \begin{beamercolorbox}[wd=\textwidth,sep=0.10cm]{block body}
    \centering\scriptsize
    \textbf{1--7:} monomio. \quad \textbf{8--12:} índice $n$. \quad
    \textbf{13--20:} conjugado y aplicaciones.
  \end{beamercolorbox}
  \listatiny
  \vspace{0.15cm}
  \begin{columns}[T,onlytextwidth]
    \begin{column}{0.47\textwidth}
      \begin{enumerate}\setlength{\itemsep}{0.16cm}
        \item $\dfrac{1}{\sqrt{7}}$
        \item $\dfrac{6}{\sqrt{3}}$
        \item $\dfrac{10}{\sqrt{5}}$
        \item $\dfrac{3}{2\sqrt{2}}$
        \item $\dfrac{\sqrt{2}}{\sqrt{5}}$
        \item $\dfrac{15}{\sqrt{10}}$
        \item $\dfrac{a}{\sqrt{b}}$
        \item $\dfrac{1}{\sqrt[3]{5}}$
        \item $\dfrac{4}{\sqrt[3]{2}}$
        \item $\dfrac{3}{\sqrt[3]{4}}$
      \end{enumerate}
    \end{column}
    \begin{column}{0.50\textwidth}
      \begin{enumerate}\setlength{\itemsep}{0.16cm}
        \setcounter{enumi}{10}
        \item $\dfrac{1}{\sqrt[4]{2}}$
        \item $\dfrac{5}{\sqrt[5]{16}}$
        \item $\dfrac{1}{\sqrt{3}+1}$
        \item $\dfrac{4}{\sqrt{6}-\sqrt{2}}$
        \item $\dfrac{2}{5-\sqrt{3}}$
        \item $\dfrac{\sqrt{3}}{\sqrt{3}-1}$
        \item $\dfrac{\sqrt{5}+1}{\sqrt{5}-1}$
        \item $\dfrac{1}{\sqrt{x}-3}$
        \item ¿Cuál es el conjugado de $7-2\sqrt{5}$?
        \item ¿Por qué racionalizar no cambia el valor?
      \end{enumerate}
    \end{column}
  \end{columns}
\end{frame}

\begin{frame}{Respuestas}{Autoevaluación de los ejercicios propuestos}
  \listatiny
  \vspace{0.15cm}
  \begin{columns}[T,onlytextwidth]
    \begin{column}{0.47\textwidth}
      \begin{enumerate}\setlength{\itemsep}{0.16cm}
        \item $\tfrac{\sqrt{7}}{7}$
        \item $2\sqrt{3}$
        \item $2\sqrt{5}$
        \item $\tfrac{3\sqrt{2}}{4}$
        \item $\tfrac{\sqrt{10}}{5}$
        \item $\tfrac{3\sqrt{10}}{2}$
        \item $\tfrac{a\sqrt{b}}{b}$
        \item $\tfrac{\sqrt[3]{25}}{5}$
        \item $2\sqrt[3]{4}$
        \item $\tfrac{3\sqrt[3]{2}}{2}$
      \end{enumerate}
    \end{column}
    \begin{column}{0.50\textwidth}
      \begin{enumerate}\setlength{\itemsep}{0.16cm}
        \setcounter{enumi}{10}
        \item $\tfrac{\sqrt[4]{8}}{2}$
        \item $\tfrac{5\sqrt[5]{2}}{2}$
        \item $\tfrac{\sqrt{3}-1}{2}$
        \item $\sqrt{6}+\sqrt{2}$
        \item $\tfrac{5+\sqrt{3}}{11}$
        \item $\tfrac{3+\sqrt{3}}{2}$
        \item $\tfrac{3+\sqrt{5}}{2}$
        \item $\tfrac{\sqrt{x}+3}{x-9}$
        \item $7+2\sqrt{5}$
        \item Porque se multiplica por una fracción igual a $1$
      \end{enumerate}
    \end{column}
  \end{columns}
\end{frame}

\begin{frame}{Referencias}
  \footnotesize
  \begin{itemize}
    \item[{[1]}] Ramos del Olmo, Á. M. (2016). \textit{Matemáticas básicas
      para el acceso a la universidad}. Ediciones Pirámide.
      \texttt{https://elibro.net/es/ereader/uniminuto/49134}
    \item[{[2]}] González Pulido, J. W., Chacón Benavides, J. A. \&
      Fonseca Correa, L. Á. (2023). \textit{Matemática básica} (1.\textsuperscript{a} ed.).
      Universidad Mariana.
      \texttt{https://elibro.net/es/ereader/uniminuto/280647}
    \item[{[3]}] Placencia Valero, J. (2008). \textit{Compendio de matemática
      básica elemental} (1.\textsuperscript{a} ed.). Editorial Tébar Flores.
      \texttt{https://elibro.net/es/ereader/uniminuto/269758}
  \end{itemize}

  \vspace{0.35cm}
  \begin{center}
    {\Large\color{uniAzul}\textbf{¿Preguntas?}}
  \end{center}
\end{frame}

\end{document}
